\section{Rapport : 12/05/2026}
\subsection{Existence d'un processus de Poisson}

\begin{important}
	Durant cette réunion, on a terminé la présentation des propriétés essentielles des processus de Poisson (voir \cite{king_poisson}). On a notamment montré l'existence de tels processus sur un espace mesurable $S$ dès lors que la mesure d'intensité peut s'écrire sous la forme
	\[
		\mu = \sum_{n \in \N} \mu_n,
	\]
	où les $\mu_n$ sont des mesures finies sans atome.
\end{important}


L'aspect constructif des processus de Poisson est particulièrement intéressant, car il permet de voir un processus ponctuel comme un ensemble aléatoire
\[
	\Pi = \{X_1,\dots,X_N\},
\]
où les variables $(X_n)_{n\in\N}$ et $N \sim \aP(\lambda)$ sont indépendantes. Cette représentation permet notamment de \og différencier les points \fg{}. En effet, dans la définition abstraite d'un processus ponctuel, rien ne garantit a priori la possibilité de choisir de manière cohérente des éléments $X\in \Pi$ pour une réalisation donnée.

Cette construction permet également d'introduire naturellement la notion de marginale. Par exemple, si
\[
	\Pi \sim \mathcal P(\lambda dx)
\]
sur $[0,1]$, alors conditionnellement à $N$, les variables $(X_i)_{1 \le i \le N}$ sont i.i.d. et suivent une loi uniforme sur $[0,1]$. La variable aléatoire
\[
	N \sim \mathcal P(\lambda)
\]
représente alors le nombre total de points du processus.

\subsection{Processus marqués}

Cette remarque est importante pour introduire les processus ponctuels marqués. On définit un processus marqué
\[
	\Pi^* = {(X,m_X) : X \in \Pi},
\]
comme un processus ponctuel à valeurs dans un espace produit $S\times M$, où $M$ désigne l'espace des marques. L'idée consiste à associer une marque à chaque point du processus tout en conservant une structure d'indépendance entre les différentes marques.

Plus précisément, si
\[
	\Pi = \{X_1,\dots,X_N\}
\]
est un processus de Poisson et si $p(x,dm)$ est un noyau de probabilité, alors les marques sont choisies de sorte que
\[
	m_{X_i} \sim p(X_i,dm),
\]
indépendamment les unes des autres conditionnellement au processus initial.

\begin{important}
	Dans ce cas, le processus marqué $\Pi^*$ est encore un processus de Poisson d'intensité
	\[
		\mu^*(dx,dm)=\mu(dx)p(x,dm).
	\]
\end{important}

Ce résultat découle directement de la formule de Campbell, qui permet de caractériser les processus de Poisson. Pour un processus ponctuel $\Pi$, on peut associer une mesure de comptage $N$ définie par
\[
	N(f)=\sum_{X\in\Pi} f(X)=\int fdN,
\]
pour toute fonction mesurable positive $f$.

\subsection{Formule de Campbell et Application}

\begin{important}
	La formule de Campbell s'écrit
	\[
		\E\left(e^{-N(f)}\right)
		=
		\exp\left(
		-\int (1-e^{-f}),d\mu
		\right).
	\]
	Réciproquement, si cette relation est satisfaite pour toute fonction mesurable positive $f$, alors le processus ponctuel considéré est un processus de Poisson d'intensité $\mu$.
\end{important}

Une application importante de cette formule concerne les processus ponctuels sur $\R_+$. Si $\Pi$ est un processus de Poisson homogène d'intensité $\lambda>0$, et si l'on ordonne ses points sous la forme
\[
	\Pi=\{X_1,\dots,X_n,\dots\},
\]
alors les temps d'attente
\[
	X_{i+1}-X_i,
	\qquad i\in\N,
\]
(avec la convention $X_0=0$) sont i.i.d. et suivent une loi exponentielle de paramètre $\lambda$.

Pour revenir à l'EDS introduite dans l'article sur les équations de fragmentation (voir \cite{frag_deaconu}), l'idée consiste à reconstruire le processus de fragmentation à partir d'un processus de Poisson ponctuel tridimensionnel
\[
	N(ds,du,dv),
\]
d'intensité $ds,du,dv$.

\begin{important}
	Le processus $(\xi_t)_{t\geq0}$ représentant la masse d'une particule typique est alors solution d'une équation différentielle stochastique à sauts de la forme
	\[
		\xi_t =
		\xi_0
		+
		\int_0^t\int_{\R_+^2}
		H(\xi_{s-},u,v),N(ds,du,dv),
	\]
	où la fonction
	\[
		H(x,u,v) =
		-u,\mathbf 1_{{u\le x}}\mathbf 1_{{v\le g^{LM}(x,u)}}
	\]
	décrit l'effet d'une fragmentation sur la masse de la particule. Cette écriture permet de représenter l'évolution du système comme une accumulation de sauts pilotés par le processus de Poisson ponctuel.
\end{important}

Cette construction est directement reliée à la chaîne de Markov introduite dans l'article. En effet, lorsque le processus est dans un état $x$, les points du processus de Poisson sont acceptés uniquement dans la région
\[
	\Gamma(x)={(u,v)\in\R_+^2 : v\le g^{LM}(x,u)}.
\]
L'aire de cette région vaut
\[
	|\Gamma(x)|
	=
	\int_0^\infty g^{LM}(x,u)du
	=
	F(x),
\]
ce qui implique que le temps d'attente avant le prochain saut suit une loi exponentielle de paramètre $F(x)$. De plus, conditionnellement au saut, la nouvelle masse est distribuée selon le noyau de fragmentation associé. On retrouve ainsi exactement la dynamique de la chaîne de Markov construite précédemment.
